\documentclass[11pt,a4paper]{article}
\usepackage[margin=0.85in]{geometry}
\usepackage{hyperref}
\title{Sector-Specific Complete-Text Clustering and Evidence-Linked Coding}
\author{}
\date{1 October 2026}
\begin{document}
\maketitle

\section{Implemented workflow}
\path{src\cdp_text_clustering\cluster_cdp_fields_by_sector.py} is the current
sector-specific refinement. It reuses the complete-field embeddings from the
pipeline below and writes a separate result directory. It does not overwrite
the original complete-versus-truncated benchmark, encoder caches, source
texts, or coding templates.

\path{src\cdp_text_clustering\benchmark_cdp_field_chunks.py} is a separate, local
pipeline. It does not overwrite the encoder competition or earlier cluster
assignments. It uses the competition's saved company-separated sample and
can select a smaller development subset without moving companies between
training and holdout.

For longitudinal initiatives, it reads
\texttt{initiative\_comment\_reported} from the currency-enriched extract.
For risks and opportunities, it independently processes
\texttt{reported\_outcome\_text} (description),
\texttt{response\_strategy\_reported}, and
\texttt{financial\_impact\_description}.
For the separate 2024 questions it reads Q7.55.2 column 9,
Q7.55.3 column 2, and Q7.55.4 column 1.
The selected CDP categories, currency, costs, and source coordinates are
metadata, not embedding input. Blank fields remain in the availability
ledger; they are not assigned zero embeddings or interpreted as absent action.

Separation is limited by the existing extract: legacy strategy columns can
include explanations of costs, and a financial-description field can
combine source cells. The pipeline preserves those fields exactly rather
than pretending to have recovered finer raw-cell distinctions.

\section{Complete-text chunking and aggregation}
Each nonblank field is partitioned into nonoverlapping chunks. Sentence
and paragraph boundaries are preferred using a multilingual punctuation
heuristic; oversized sentences use a token-checked character/whitespace
fallback. This is not a linguistically validated sentence segmenter.
All original characters, including whitespace, occur exactly once.
Each chunk stores its source field, chunk ID, exact text, start/end
coordinates, token count, and aggregation weight.

Coordinates are zero-based Python Unicode character offsets with an
exclusive end. They refer to the exact \emph{extracted source field},
not a byte offset, translated answer, or original workbook cell.
The actual encoder tokenizer counts the model-specific prompt and special
tokens. A chunk exceeding the budget causes an explicit failure rather
than silent truncation.

The default uses locally downloaded MiniLM at its distributed
128-token limit. E5, BGE-M3, and Qwen3 are selectable using their pinned
competition revisions and model-specific prefixes. The script is offline
and requires the chosen weights to exist locally.

Chunk vectors are pooled with weights equal to their content-token counts
(at least one), then the weighted mean is normalized. There is exactly
one embedding per response/source-field pair. Fields are never pooled
together, and chunks never become independent observations in clustering.
Nonoverlap prevents double-counting source tokens through overlapping
windows. Boundary tokenization can still differ from whole-text tokenization.
Mean pooling may blur multiple themes; it is a benchmark representation,
not a substitute for multi-label coding.

\section{Original matched comparison (preserved baseline)}
Within each dataset and source field, the script compares complete-text
weighted pooling against truncation of that \emph{same separated field}.
Both conditions use the same encoder, records, company split, cluster
count, and five initialization seeds. This isolates text handling more
cleanly than comparing new fields with the old concatenated narrative.

The benchmark fits one row per response/field: twelve clusters, except six
for Q7.55.4. If there are insufficient distinct training fields or eligible
holdouts, the condition is explicitly recorded as skipped. Exact normalized
field texts appearing in training are excluded from holdout diagnostics
and listed in \texttt{holdout\_eligibility.csv}. Such fields remain in
the source registry and coding workflow.

The outputs contain silhouettes in both the complete and truncated
representation spaces, seed stability, and adjusted Rand agreement between
text-handling conditions. Neither agreement nor separation establishes
human accuracy. A changed cluster assignment can be an improvement or an
error. The held-out sample is a development diagnostic, not an untouched
final scientific test.

\section{Sector-specific refinement (current clustering)}
The current stage partitions the data by \emph{dataset, source field, and
reported CDP primary industry} before fitting any model. Thus, a cluster
cannot combine sectors, risks with initiatives, or comments with financial
explanations. CDP primary industry is the broad, comparable classification
available across the source formats. The finer reported primary sector is
preserved separately, not silently mixed with the broader level. Industry is
joined from the original source record and checked against company, year,
and exact field text; it is not guessed from keywords or company names.
Missing classifications remain explicitly unclustered.

\subsection{More specific representation and cluster counts}
Within each eligible sector/field cohort, a training-fitted TF--IDF
representation of words and two-word phrases complements the existing
normalized complete-field embedding. The blocks are scaled by
$\sqrt{0.7}$ and $\sqrt{0.3}$ and concatenated, then each row is normalized.
This retains multilingual semantic information while giving distinctive
mechanisms such as LED lighting, compressed air, solar panels, and heat
recovery more influence than in a semantic-only representation.
Generic reporting and financial boilerplate is removed from the lexical
vocabulary, but the semantic input and canonical source text are unchanged.
There is still one model observation per response/source field, not per
chunk or per suggested code.

Neither the vocabulary nor the clustering model is fitted on holdout rows.
Normalized duplicate training texts contribute only one fitting observation.
Candidate cluster counts range from two to sixteen, bounded by the available
unique training fields. Every candidate is fitted with seeds 42, 123, and
2026, using ten K-means initializations per seed. A candidate is supported
only if \emph{every seed} produces at least three unique training fields and
two distinct companies in every cluster. A cohort needs at least six unique
training fields and four training companies before a search is attempted.

The search uses mean training cosine silhouette. Among supported positive
scores within 0.02 of the best score, the largest cluster count is selected.
This explicitly prefers finer, similarly separated partitions rather than
fixing twelve clusters regardless of sector size. The best-scoring seed for
that count supplies the saved model; adjusted Rand agreement across the
three seeds is reported as a stability diagnostic. These are exploratory
settings, not evidence that sixteen clusters or a 0.02 tolerance is optimal.

Sparse sectors, holdout-only sectors, and cohorts without a supported
positive-separation partition are recorded as unclustered with a reason.
They are not pooled into an ``other'' sector, forced into singletons, or
silently discarded. Blank fields also remain in the assignment ledger.
Eligible holdouts are predicted by their own sector's training model only.
Exact normalized holdout duplicates of training text in the same dataset
and source field are excluded from holdout diagnostics, even when the
duplicate is in another sector. They can still receive a model assignment.

\subsection{Specific labels and evidence remain distinct}
Cluster names contain the sector and up to five distinctive training
terms, scored by within-cluster versus other-cluster mean TF--IDF.
Representative examples are the three nearest unique training fields.
Neither holdout text nor holdout labels is used to name clusters.
If a training cohort has no lexical tokens, it explicitly uses semantic
features only and exports a numbered semantic descriptor.

The refined 67-code vocabulary distinguishes, for example, lighting
efficiency, IT efficiency, fleet electrification, renewable procurement,
solar generation, heat recovery, and gas flaring reduction.
The 182 development initiative reviews supply 200 exact evidence-linked
code suggestions as a separate overlay on matching full-sample fields.
Their canonical initiative names and codes are visible in the assignments;
per-cluster code counts use only reviewed unique training fields.
These suggestions are \emph{not} clustering features, constraints, or
human-accepted ground truth. In particular, a code observed on one reviewed
member is not automatically assigned to every other cluster member.
Unreviewed fields have descriptive model-cluster names, not fabricated
canonical code annotations.

For more precise initiative grouping, a parallel evidence-guided layer
groups the reviewed fields by \emph{sector plus canonical initiative code}.
These groups can overlap for multi-mechanism responses and can link
longitudinal and 2024 initiative comments through their shared code.
IT efficiency and paper reduction therefore remain distinct even when
K-means places them in the same broad office-related topic. The assignment
table contains \texttt{specific\_initiative\_group\_ids}; these are separate
from the learned \texttt{cluster\_id}. Small evidence groups are retained
with support flags, not misrepresented as fitted models. Missing-sector and
insufficient-mechanism reviews are not assigned guessed specific groups.
This layer uses the pending development reviews, including reviewed holdout
fields, and is not an independent accuracy evaluation.

One model cluster per field does not imply one action per response.
Multi-label coding remains separate; planned and implemented actions are
not distinguished by cluster membership alone. Assignment cosine and the
distance margin between the nearest two centroids are descriptive measures,
not probabilities. Negative or undefined holdout silhouettes, sparse-sector
exclusions, and ambiguous examples must be considered before substantive use.

\section{Evidence-linked multi-label coding}
The draft English codebook is
\path{src/cdp_text_clustering/cdp_field_codebook.json}. Separate roles cover actions,
risks, opportunities, financial consequences, funding mechanisms, and
barriers. These definitions require human refinement and are not a
validated taxonomy.

In the original benchmark, for each chunk the model retrieves the two closest applicable code
definitions and writes \texttt{coding\_candidates.jsonl}. These are
\emph{unreviewed retrieval suggestions}, with cosine similarities rather
than correctness probabilities. Their evidence spans initially cover
whole chunks; the reviewer should narrow them to the supporting words.
Negation, planned versus completed status, contradictions, and insufficient
information still require human judgment. A high similarity is not evidence
that an action occurred.

\texttt{review\_fields.csv} provides source context without code suggestions.
Externally supplied text is formula-escaped for spreadsheet display;
\texttt{fields.jsonl} remains the exact canonical text for offset checks.
\texttt{coding\_template.csv} initially has one pending row per available field.
Repeat rows to code multiple themes or supply several evidence spans.
Fill in a code ID, \texttt{evidence\_start}, \texttt{evidence\_end},
\texttt{evidence\_text}, and \texttt{reviewer\_id}; use review status
\texttt{accepted}, \texttt{rejected}, or \texttt{pending}.
Existing templates are never overwritten by reruns.

The development review stage,
\path{src\cdp_text_clustering\review_cdp_initiatives.py}, replaces initiative
retrieval suggestions with the evidence-linked review manifest and
populates blank pending template rows, using multiple rows where necessary.
Nonblank annotation work is preserved. Original development candidates,
template, and codebook are archived as \texttt{*.original.*}.
Sixteen initiative comments are marked as insufficient information rather
than assigned an unsupported mechanism; 747 fields outside the two initiative
datasets retain their original suggestions.
Both training and holdout examples informed this development vocabulary,
so a separate untouched sample is needed to assess coding accuracy.

Validation rejects unknown/inapplicable codes, invalid source identities,
invalid offsets, mismatched excerpts, duplicate entries, and accepted rows
without a reviewer. It reports pending/rejected counts explicitly.
Accepted evidence is exported in long form and aggregated into per-response
code sets \emph{separately for each reviewer}; disagreements are not
silently adjudicated. No human labels are manufactured by the benchmark.
The validator verifies evidence linkage, not whether a human interpretation
is substantively correct. Empty accepted output means no accepted coding,
not an absence of climate action.

\section{Execution and outputs}
\subsection{Preserved source benchmark}
The completed source MiniLM run is saved under
\path{data/outputs/cdp_text_clustering/cdp_field_chunks_20261001/}.
It processes all 3,579 responses in the matched competition sample:
5,615 nonblank fields and 364 blank fields. The 8,729 chunks retain
all 2,828,733 source characters, with no chunk exceeding 128 input tokens.
There are 17,458 unreviewed top-two code suggestions, not 17,458 validated
actions or independent observations.

All ten dataset/field comparisons completed, producing 100 fits across
two text-handling conditions and five seeds. Two holdout fields were
excluded because their exact normalized text occurred in training.
Complete-versus-truncated partition agreement ranges from approximately
0.300 to 0.564 across these comparisons. That difference demonstrates
sensitivity to text handling, not that changed assignments are correct.
The single-chunk response vectors match the prefix comparator up to
floating-point rounding. No accepted human labels were generated.
The field-aware run has been executed for MiniLM; support for the other
encoders does not imply their chunked experiments have already run.

Run with the configured project Python interpreter containing the existing
NumPy, SciPy, scikit-learn, pandas, joblib, and threadpoolctl dependencies.
The original encoder pipeline additionally requires its installed encoder
dependencies and local model weights. The source benchmark processes up to
600 sampled responses per dataset. To rerun it with the expanded codebook,
choose a new directory rather than overwriting the recorded 36-code run:
\begin{verbatim}
python -m src.cdp_text_clustering.benchmark_cdp_field_chunks `
  --output data\outputs\cdp_text_clustering\cdp_field_chunks_67code_baseline
\end{verbatim}
A bounded development run across all six datasets:
\begin{verbatim}
python -m src.cdp_text_clustering.benchmark_cdp_field_chunks `
  --per-dataset 100 `
  --output data\outputs\cdp_text_clustering\cdp_field_chunks_minilm_development_v2
\end{verbatim}
The examples use PowerShell backticks for line continuation.
Use \texttt{--model e5}, \texttt{bge\_m3}, or \texttt{qwen3} with a
different output directory to compare another locally available encoder.
Only the selected sample is processed, not all 314,298 longitudinal records.

\subsection{Current sector-specific output}
The full-sample sector refinement produces \textbf{236 learned clusters}
from 84 fitted sector/field cohorts; 50 of those clusters belong to the two
initiative datasets. No learned cluster crosses a sector boundary.
There are 127 observed nonblank, known-sector cohorts in total.
Coverage of the complete source registry is:

\begin{center}
\begin{tabular}{lr}
Assignment status & Fields \\
\hline
Assigned to a fitted sector cluster & 4,320 \\
Missing reported sector/industry & 862 \\
No supported multi-cluster partition & 365 \\
Insufficient unique training/company support & 68 \\
Blank source field & 364 \\
\hline
Total source fields & 5,979
\end{tabular}
\end{center}

Thus, 4,320 of the 5,615 nonblank fields receive model assignments.
Of the 84 fitted cohorts, 54 have a defined holdout silhouette; the remaining
30 have insufficient holdout observations or distinct predicted clusters
for that metric. Undefined metrics are not reported as zero.
The evidence-guided initiative layer contains \textbf{91 sector--mechanism
groups}, with 171 memberships across 153 distinct reviewed fields.
Of those groups, 35 contain multiple fields and 56 contain only one reviewed
field; the latter are coding groups awaiting more evidence, not reliable
statistical subclusters. The other reviewed fields lack a known sector or a
specific supported mechanism. All 182 initiative reviews and all 200 code
suggestions remain available in the evidence overlay.

Run the current stage against the already completed full source benchmark:
\begin{verbatim}
python -m src.cdp_text_clustering.cluster_cdp_fields_by_sector `
  --source data\outputs\cdp_text_clustering\cdp_field_chunks_20261001 `
  --output data\outputs\cdp_text_clustering\cdp_field_chunks_sector_specific_20261001
\end{verbatim}
No encoder download or repeated embedding inference is needed. The output
fingerprint records input hashes, the sector join, codebook and review
manifest, feature weights, seeds, and granularity settings. Different inputs
or settings require a different output directory.

\texttt{cluster\_assignments.csv} contains every original field with its
sector, cluster or explicit unclustered status, source text, and available
reviewed initiative names. \texttt{cluster\_summary.csv} contains specific
cluster names, training/holdout sizes, company support, and reviewed
training-code counts. \texttt{cohort\_metrics.csv} records sector-level
eligibility, selected granularity, stability, and holdout diagnostics;
\texttt{granularity\_search.csv} retains every candidate count and seed,
including rejected small-cluster partitions.
\texttt{representative\_fields.csv} provides training examples for review.
For the most specific reviewed initiative groups, use
\texttt{initiative\_specific\_groups.csv} and
\texttt{initiative\_specific\_group\_summary.csv}. These preserve the
sector, canonical action name, exact evidence, support counts, and crosswalk
to the learned model cluster. Do not conflate their overlapping memberships
with independent response counts.
The \texttt{models} directory stores each fitted vectorizer and K-means
model with its exact fitting field IDs.
\texttt{reviewed\_initiative\_evidence.jsonl} retains the 200 pending
evidence suggestions and offsets into the unchanged source registry.
\texttt{completion.json} records coverage and limitations after the stage
has completed; it is not a human-validation certificate.

After editing a copy of the coding template:
\begin{verbatim}
python -m src.cdp_text_clustering.benchmark_cdp_field_chunks `
  --stage validate `
  --output data\outputs\cdp_text_clustering\cdp_field_chunks_minilm_development `
  --annotations path\to\reviewed_codes.csv
\end{verbatim}
The annotation input can alternatively be JSONL. For independent validation,
hide AI suggestions and have coders read the complete source fields first.

Key artifacts include \texttt{fields.jsonl}, \texttt{chunks.jsonl},
\texttt{coverage\_audit.csv}, \texttt{response\_field\_embeddings.npy},
its row manifest, paired benchmark metrics, saved models, assignment tables,
the draft codebook, and human-review files.
Successful runs write \texttt{completion.json} only after all stages finish.
Token-budget and full-character-coverage checks run before inference.
Embedding blocks are checkpointed and reused; different inputs, codebooks,
model revisions, or settings require a new output directory.

Focused unit tests (from the repository root):
\begin{verbatim}
Push-Location src\data_engine
python -m unittest test_cdp_field_chunking `
  test_review_cdp_initiatives test_cluster_cdp_fields_by_sector -v
Pop-Location
\end{verbatim}
\end{document}
